\section{Putting a Rose Back on the Cross}

Cologero has been very generous in allowing many of us to write on Gornahoor.

For that reason, I've sometimes repaid him by not posting anything, when I was sure I had nothing important or clear to say.

I am ending the series on Iamblichus; the reader will have to write his own recapitulation, or Decad. I am not certain where I will go from here, but have purposed not to write ``for the sake of writing''.

There are a couple of thoughts, I think, worth re-visiting, concerning Christian esotericism.

1. Many of the texts seem to focus on obtaining the virtues. Is this because Christianity is lost in the ``lesser mysteries'' of obtaining back primordial man? I am inclined not to think so. Whereas much neo-pagan initiation seems to center around an ``experience'' or a realization, Christian practice (or those influenced by it) invariably focus on acquiring concrete moral virtues. These concrete moral virtues are meant to lead to higher angelic states, symbolized in the theological virtues.

2. This contrast is even seen with one of the grandest Stoics, Epictetus.\footnote{\url{http://en.wikipedia.org/wiki/Epictetus}} There is nothing but good to say of this man, but when all is said and done, he is not the same figure as the Apostle Paul: he does not present us with the theological virtues. This is not his fault, nor even a fault. It is simply an observation.

3. In the Romance of the Rose\footnote{\url{http://www.legends-and-myths.com/40_2.cfm?i=1001-legends-myths-initiation-romance-rose}}, we are presented with a garden whose wall is guarded by several figures representing Envy, Sorrow, Poverty, etc., which the inquirer must ``pass'' or eschew in order to enter into the garden. Or do these disfigurements (in another sense) actually function as guardians? In any case, it is hard to miss the parallel with the Seven Deadly Sins. Again, Christianity is definitely aimed at conquering manifested sins in a manifested way; it is a religion of Incarnation.

4. While it thus (granted) appeals to those at the lowest level intellectually, it maintains higher links capable of being articulated within the most subtle and refined cosmologies or metaphysics. I have already in the past endeavored to link Steiner's work on the Seven Steps of Christianity\footnote{\url{https://gornahoor.net/?p=2880}} with a passage from Saint Peter, and indirectly, the Sermon on the Mount. Granted (again), this is cosmology or hermeticism, which as Guenon points out, is not a complete metaphysic, and we follow Guenon on that point. Part of the reason Christianity is hard to define is that it can fit various metaphysics.

5. Where does this leave the esoteric Christian, or the Christian who wants to be faithful to his tradition (and thereby, to God), but go beyond the low and easily distorted spirituality which passes for currency in today's world of cheap grace? Guenon even argues that the Rosicrucian tradition (a hermetic one, and so incomplete, devoted to the lesser mysteries) has passed into desuetude. Charles Upton makes the same claim regarding the neopagan mystery school of Iamblichus. A great deal has been ``lost''. But is it really lost? Is it really embedded in Time in such a fashion that it is incapable of rekindling embers and starting a new fire? We think not.

6. This leaves the task of our day; as Cologero put it, we should not sit around waiting for George to do it. On the contrary, we are not wine-snobs and connoisseurs waiting to be fed -- our job is to begin to prepare the meal, if we want to eat.

7. We can always start with attention to our thought and our motivations and thereby, to the deepest wellsprings of our actions, whether spiritual or physical. A resurgent body of Christian esotericism no longer has to contend with infra-Church bodies of Inquisitors, political machinations of suspicious kings, or a host of other difficulties encountered during the Middle Ages (which, all in all, were a better time). We also have more information and resources, of a kind, at our disposal. In this area, at least, much will be required of us. Then, perhaps as Cologero has said, if we prepare these lesser things, and do our homework, the greater mystery of what spark will rekindle the flames will be revealed. In any case, it is undoubtedly still alive -- what would the ``Wise Men'' of today do to experience the birth of the Christ-child, but look to the sign of the Eagle and the Cross? Then, perhaps, we will regain the Rose.

8. As a personal postscript, I do think Charles Williams must have had some knowledge, however fragmentary, of what we seek here --- his writings center around too many hermetic themes. Even his interest in Dante is remarkable. Coupled with his early (relatively) death, he is a ``person of interest'' to me.


\flrightit{Posted on 2013-04-19 by Logres}

\begin{center}* * *\end{center}

\begin{footnotesize}
\begin{sffamily}
\texttt{George on 2013-04-19 at 22:43 said:}

Thank you Logres. I find your posts very interesting and helpful.

\hfill

\texttt{J. Alexander Maximilian on 2013-04-20 at 05:28 said:}

I too really enjoyed your writings, and hope to see your writings again soon.

God bless.

\hfill

\texttt{David on 2013-04-20 at 09:23 said:}

Very interesting. As you say, we must be ever vigilant and responsible, just as the parables of Mt 25. This particular sentence was very inspiring : «Again, Christianity is definitely aimed at conquering manifested sins in a manifested way; it is a religion of Incarnation». Thank you for your writings.

\hfill

\texttt{George on 2013-04-20 at 16:47 said:}

Would it be possible to create a ``lesson plan'' for studying Christian esoterism, with levels of reading from beginning through advanced? I think this could help some uneducated folks (like myself).

\hfill

\texttt{George on 2013-04-20 at 19:46 said:}

Would any recommend ``Guenonian Esoterism And Christian Mystery''?

\hfill

\texttt{J. Alexander Maximilian on 2013-04-20 at 19:54 said:}

George, I would highly recommend ``Meditations on the Tarot: A Journey Into Christian Hermeticism'' a most excellent book for beginner and master alike. Also study Arthuriana deeply. In a way at this time the Christian Hermetic mystic has to forge a path, luckily I have experience in the Golden Dawn and OTO. They helped greatly in my quest however I would tell the neophyte to avoid those now and focus on Christian mysteries.

Look inti Angelology, and just do internet searches for Christian Hermarica, etc. But really ``Meditations'' is an excellent start.

\hfill

\texttt{Logres on 2013-04-20 at 22:04 said:}

I can't comment on the book, but Jean Borella seems fairly sound (I have only read one of his works). He is of interest, if for no other reason, because he is a theologian in the Church who seems to have taken Guenon very seriously. Others, like F. Seraphim Rose, freely admit their debt, but ``moved on'' afterwards. Any Christian who wants to grow in their faith, and is interested in traditional teaching should probably read this or books like them. George Heart's work on Gnosis is very good -- Cologero recommended that one.

\url{http://www.amazon.com/Christianity-Dogmatic-Gnostic-Vivifying-Knowledge/dp/1412062284}

Robin Amis has done some work on Christian esotericism.

\url{http://praxisresearch.net/}

He even has a website. I tried to put together such a primer, but I don't trust my judgement enough yet. There has been a lot of criticism (see the Forum) directed at people like James Cutsinger and others who (at first glance) seem to be paralleling ``the Faith'' with some kind of overarching Tradition with a capital T. Since I believe in a God simultaneously man and yet God, three and yet One, it doesn't bother me to think that there is salvation outside of ``the Faith'': the Faith remains the Faith.

\url{http://en.wikipedia.org/wiki/James_Cutsinger}

All Tradition is under attack today -- that's why they call it the Kali Yuga. We see how ``traditional'' Islam is now having suicide bombers carry out their attacks on civilians. So everything that can be subverted, will be. What might be helpful is to aid one another in studying various branches of Tradition that individuals can ``see'' where to fit within the Faith, without dilution or compromise. This is a difficult task; but what else does ``every thought captive to Christ'' entail? Mihai is working on this, by the way. Father Thaddeus' little book on Thoughts, or ``Unseen Warfare'' by Scupoli are good places to start. I think we will have to go back to spiritual fundamentals, in order to get anywhere recovering (for example) the Rosicrucian initiation (which supposedly derived from Apostle Mark's conversion of some Egyptian theurgists). If the research is done with enough right reason and honesty, and spiritual fundamentals like humility and vigilance are practiced, a primer could be done. Perhaps that should be attempted.

\hfill

\texttt{Logres on 2013-04-20 at 22:08 said:}

Agree wholeheartedly with the advice given by J Maximillian. I almost went down a dangerous path, had Cologero not recommended Meditations. Definitely read this, and see Cologero's website for excerpts. The beauty of Christianity is that rapid progress can be made through humility and charity: it's important to remember that our interest in Tradition be incidental, and entirely subordinate to God. Otherwise, we risk falling into ``magic'', the opposite of Tassawuf (as the Sufis would say -- credit to Charles Upton on these observations). At least for the Christian, who is without excuse in these matters.

And thanks to everyone for commenting.

\hfill

\texttt{Cassiodorus on 2013-04-21 at 22:57 said:}

Yes. But, I must say that it was slow going for me to absorb.

I would also recommend Wolfgang Smith's ``Christian Gnosis'' who offers some interesting views on the relationship of Trinitarian theology and Perennialist metaphysics.

A much easier introductory work that I thought was quite good would be Richard Smoley's ``Inner Christianity: A Guide to the Esoteric Tradition''.

\hfill

\texttt{Jason-Adam on 2013-04-22 at 12:04 said:}

What is your opinion on Arthur Edward Waite ?

Borella is one of the best. Tomberg is great as well. Also I can recommend Eliphas Levi, a real Catholic Cabalist.

Three Ages of the Interior Life by Father Gerrigou-Lagrange is a good primer on mystical theology.

Try to read up on antique and mediaeval philosophy first though -- Plato, Aristotle, Plotinius, St Dionysius, St Augustine, St Boethius, St Thomas Aquinas, St Bonaventura, Meister Eckhart, St John of the Cross etc etc etc

\end{sffamily}
\end{footnotesize}

